\documentclass[11pt]{article}
\usepackage[a4paper,margin=26mm]{geometry}
\usepackage[T1]{fontenc}
\usepackage{lmodern}
\usepackage{amsmath,amssymb,amsthm}
\usepackage{microtype}
\usepackage{booktabs}
\usepackage{array}
\usepackage{xurl}
\usepackage{hyperref}
\hypersetup{colorlinks=true,linkcolor=blue,urlcolor=blue,citecolor=blue,
 pdftitle={Tensor products of Mycielskians of odd cycles with odd wheels are not word-representable},
 pdfauthor={Horace Dong}}
\newtheorem{theorem}{Theorem}
\newtheorem{lemma}{Lemma}[section]
\newtheorem{proposition}[lemma]{Proposition}
\theoremstyle{remark}
\newtheorem{remark}[lemma]{Remark}
\newcommand{\Z}{\mathbb Z}
\DeclareUrlCommand\code{\urlstyle{tt}\def\UrlBreaks{\do\_\do\/\do\.}\def\UrlBigBreaks{}}
% Repository version 1.9.0 (user decision 2026-09-28T19:22:59Z: released together with version 2 of Paper 4).
\newcommand{\repoversion}{1.9.0}
\setlength{\emergencystretch}{1.5em}

\title{Tensor products of Mycielskians of odd cycles\\ with odd wheels are not
word-representable}
\author{Horace Dong\\ \small Independent researcher\\
\small\texttt{maxwelldhx@gmail.com}\\
\small ORCID: \href{https://orcid.org/0009-0002-7554-7548}{0009-0002-7554-7548}}
\date{September 2026}
\begin{document}
\maketitle

\begin{abstract}
A graph is word-representable if there is a word over its vertex set,
containing every vertex, in which two distinct letters alternate exactly when
the corresponding vertices are adjacent. Alshammari asked whether the tensor products
$\mu(C_{2n+1})\times W_{2m+1}$ and $\mu'(C_{2n+1})\times W_{2m+1}$ are
non-word-representable for all $n\ge1$ and $m\ge2$, where $\mu$ and $\mu'$
denote the Mycielskian and the extended Mycielskian and $W_{2m+1}$ is the
wheel on an odd cycle. We show that the answer is yes. When $n\ge m$, each
product contains an induced copy of its first factor; when $n\le m$, both
contain an induced copy of $\mu(C_{2m+1})$. These graphs are not
word-representable by theorems of Hameed and of Kitaev and Pyatkin.
The result, the non-word-representability of $\mu(C_{2t+1})$ and
$\mu'(C_{2t+1})$ on which it relies, and the other facts its proof uses are
formally verified in Lean~4. The proof was found by AI agents
working under the author's direction; their role is described at the end of
the paper.
\end{abstract}

\section{Introduction}
All graphs are finite and simple. Two distinct letters $x$ and $y$
\emph{alternate} in a word $w$ if deleting from $w$ all letters other than
$x$ and $y$ leaves a word $xyxy\cdots$ or $yxyx\cdots$, of even or odd
length. A graph $G$ is \emph{word-representable} if there is a word $w$ over
the alphabet $V(G)$ in which every vertex occurs and, for all distinct
vertices $x$ and $y$, the letters $x$ and $y$ alternate in $w$ if and only if
$xy$ is an edge~\cite{kl,alsh}. Halld\'orsson, Kitaev and Pyatkin~\cite{hkp}
proved that a graph is word-representable if and only if it has a
semi-transitive orientation.

The \emph{tensor product} $G\times H$ of graphs $G$ and $H$ has vertex set
$V(G)\times V(H)$, and $(a,b)$ is adjacent to $(a',b')$ if and only if $aa'$
is an edge of $G$ and $bb'$ is an edge of $H$. Kitaev and Lozin asked whether
anything definite can be said about the word-representability of tensor
products of two word-representable graphs, of a word-representable and a
non-word-representable graph, or of two non-word-representable graphs
\cite[Problem~7.2.5]{kl}, as quoted in~\cite[Problem~1.1]{alsh}. Choi, Kim and
Kim~\cite{ckk} showed that the tensor product of two graphs with
semi-transitive orientations need not have one. Alshammari~\cite{alsh}
studied tensor products of graphs containing wheels, complete graphs and
(extended) Mycielskians of cycles, and posed four problems. The first of them
is the following~\cite[Problem~5.1]{alsh}.

\begin{quote}
Is it true that $\mu_{2n+1}\times W_{2m+1}$ and $\mu'_{2n+1}\times
W_{2m+1}$ is always non-word-representable for $n\ge1$ and $m\ge2$?
\end{quote}

Here, as in~\cite{alsh}, $W_k$ is the wheel obtained from the cycle $C_k$ by
adding a vertex adjacent to all its vertices, $\mu_k=\mu(C_k)$ is the
Mycielskian of $C_k$ and $\mu'_k=\mu'(C_k)$ its extended Mycielskian
(Section~\ref{sec:defs}). We answer the question in the affirmative.

\begin{theorem}\label{thm:main}
For all integers $n\ge1$ and $m\ge2$, the graphs $\mu_{2n+1}\times W_{2m+1}$
and $\mu'_{2n+1}\times W_{2m+1}$ are not word-representable.
\end{theorem}

The proof (Section~\ref{sec:proof}) finds a non-word-representable induced
subgraph: for $n\ge m$, a copy of the first factor, embedded through a
homomorphism into $W_{2m+1}$, as in the proof
of~\cite[Theorem~3.6]{alsh}; for $n\le m$, a copy of $\mu_{2m+1}$, by an
embedding that adapts the one in the proof of~\cite[Theorem~4.3]{alsh}. That
these subgraphs are not word-representable is due to Hameed~\cite{hameed} and
to Kitaev and Pyatkin~\cite{kp}. We claim nothing about the other problems
of~\cite{alsh}; in particular, whether the tensor product of two
non-word-representable graphs can be word-representable~\cite[Problem~5.4]{alsh}
remains open.

\paragraph{Prior work.}
We found no earlier answer to~\cite[Problem~5.1]{alsh}. We searched on
27 September 2026 between 18:06 and 19:02 UTC: the arXiv (its API, for the
topic and the author); the works citing~\cite{alsh} in OpenAlex (there were
none); MathDB; MathOverflow and Mathematics Stack Exchange, through the Stack
Exchange API; GitHub (code, commits, issues, pull requests and
repositories), including the repositories \texttt{trureturing}~\cite{tru}
and \texttt{SCOPE2026}~\cite{scope} of AI-assisted work on open problems;
Wikipedia; and web search. A second AI system repeated a smaller search later
that day.
A final check between 23:38 UTC on 28 September and 00:17 UTC on
29 September 2026 (the arXiv, including new versions of~\cite{alsh}; the works
citing it in Semantic Scholar and OpenAlex; MathOverflow and Mathematics Stack
Exchange; GitHub; MathDB; and web search) found nothing new. We did not search Google
Scholar, Scopus or Web of Science, and we did not consult the book~\cite{kl}
itself. This reports the coverage of our searches, not a guarantee of
priority.

\section{Definitions and three known facts}\label{sec:defs}
For $k\ge3$ the cycle $C_k$ has vertex set $\Z_k$, with $i$ adjacent to
$i\pm1$. The wheel $W_k$ has the \emph{rim} vertices $r_i$, $i\in\Z_k$,
which form a copy of $C_k$, and the \emph{hub} $h$, adjacent to every rim
vertex. The Mycielskian $\mu_k$ has the \emph{top} vertices $v_i$, the
\emph{shadows} $u_i$ ($i\in\Z_k$) and the \emph{root} $w$; its edges are
$v_iv_{i+1}$, $u_iv_{i+1}$, $u_iv_{i-1}$ and $wu_i$ for $i\in\Z_k$. The
extended Mycielskian $\mu'_k$ has the same vertices and the edges $v_iv_{i+1}$,
$u_iv_j$ for all $j\ne i$, and $wu_i$~\cite{alsh,kp}. In both graphs the
shadows are pairwise non-adjacent, and the root is adjacent to no top vertex.

\begin{lemma}[{well known; see~\cite{kl}}]\label{lem:hered}
An induced subgraph of a word-representable graph is word-representable.
\end{lemma}
\begin{proof}
If $w$ represents $G$ and $S\subseteq V(G)$, delete from $w$ all letters
outside $S$. Every letter of $S$ still occurs, and whether two letters of
$S$ alternate depends only on the subword formed by these two letters, which
does not change; so the new word represents the subgraph induced by $S$.
\end{proof}

\begin{lemma}\label{lem:graphhom}
If $f\colon G\to H$ is a homomorphism of graphs, then $x\mapsto(x,f(x))$ is an
isomorphism of $G$ onto an induced subgraph of $G\times H$.
\end{lemma}
\begin{proof}
The map is injective. For vertices $x$ and $y$ of $G$, the pair
$(x,f(x))$, $(y,f(y))$ is an edge of $G\times H$ if and only if $xy$ is an
edge of $G$ and $f(x)f(y)$ is an edge of $H$; as $f$ is a homomorphism, this
holds if and only if $xy$ is an edge of $G$.
\end{proof}

\begin{theorem}[{Hameed~\cite{hameed}; Kitaev and Pyatkin~\cite[Theorem~4 and Remark~5]{kp}}]\label{thm:kp}
For every $t\ge1$, the graphs $\mu_{2t+1}$ and $\mu'_{2t+1}$ are not
word-representable.
\end{theorem}

Hameed~\cite{hameed} proved that $\mu'_{2t+1}$ has no semi-transitive
orientation; Kitaev and Pyatkin~\cite[Theorem~4]{kp} proved the same for
$\mu_{2t+1}$, and observed that their proof also applies to
$\mu'_{2t+1}$~\cite[Remark~5]{kp}. By~\cite{hkp}, these graphs are
therefore not word-representable. We use Theorem~\ref{thm:kp} only for
$t\ge2$.

\section{Proof of Theorem~\ref{thm:main}}\label{sec:proof}
For odd integers $K\ge k\ge3$ define the \emph{fold}
$\varphi_{K,k}\colon\Z_K\to\Z_k$ by $\varphi_{K,k}(i)=i$ for $0\le i\le k-1$,
and, for $k\le i\le K-1$, by $\varphi_{K,k}(i)=k-2$ if $i-k$ is even and
$\varphi_{K,k}(i)=k-1$ if $i-k$ is odd. For $K=k$ it is the identity. For
$K>k$ it is a homomorphism $C_K\to C_k$: the images of $i$ and $i+1$ are
adjacent for $0\le i\le K-2$, and $\varphi_{K,k}(K-1)=k-1$ is adjacent to
$\varphi_{K,k}(0)=0$, as $K-1-k$ is odd. In particular
$\varphi_{K,k}(i)\ne\varphi_{K,k}(i\pm1)$ for every $i$.

\begin{lemma}\label{lem:hom}
Let $k\ge K\ge5$ be odd. There are homomorphisms $\mu_k\to W_K$ and
$\mu'_k\to W_K$.
\end{lemma}
\begin{proof}
Put $\psi=\varphi_{k,K}$. Map $v_i$ and $u_i$ to $r_{\psi(i)}$, and $w$ to
$h$. The edges $v_iv_{i+1}$ and $u_iv_{i\pm1}$ of $\mu_k$ go to the rim edges
$r_{\psi(i)}r_{\psi(i+1)}$ and $r_{\psi(i)}r_{\psi(i\pm1)}$, and $wu_i$ goes to
the edge $hr_{\psi(i)}$. For $\mu'_k$, map $v_i$ to $r_{\psi(i)}$, every $u_i$
to $h$, and $w$ to $r_0$: the edges $v_iv_{i+1}$ go to rim edges, and the
edges $u_iv_j$ and $wu_i$ to edges at the hub.
\end{proof}

\begin{lemma}\label{lem:emb}
Let $K\ge k\ge3$ be odd, with $K\ge5$, and let $G$ be $\mu_k$ or $\mu'_k$.
Then $\mu_K$ is isomorphic to an induced subgraph of $G\times W_K$.
\end{lemma}
\begin{proof}
Put $\varphi=\varphi_{K,k}$ and define $\iota$ on the vertices of $\mu_K$ by
\[
 \iota(v_i)=(v_{\varphi(i)},r_i),\qquad \iota(u_i)=(u_{\varphi(i)},r_i),\qquad
 \iota(w)=(w,h).
\]
It is injective: the second coordinate determines $i$, or shows that the
vertex is $w$, and the first coordinate tells top vertices from shadows.

Edges go to edges. For $v_iv_{i+1}$: the top vertices $v_{\varphi(i)}$ and
$v_{\varphi(i+1)}$ are adjacent in $G$, as $\varphi$ is a homomorphism, and
$r_ir_{i+1}$ is an edge. For $u_iv_{i\pm1}$: in $\mu_k$ the shadow
$u_{\varphi(i)}$ is adjacent to $v_{\varphi(i\pm1)}$ because $\varphi(i)$ and
$\varphi(i\pm1)$ are adjacent in $C_k$, and in $\mu'_k$ because they are
different; and $r_ir_{i\pm1}$ is an edge. For $wu_i$: $w$ is adjacent to
$u_{\varphi(i)}$ in $G$, and $h$ to $r_i$.

Non-edges go to non-edges. The non-edges of $\mu_K$ are $v_iv_j$ with
$j\notin\{i-1,i,i+1\}$, $u_iv_j$ with $j\notin\{i-1,i+1\}$ (including
$j=i$), $u_iu_j$ with $j\ne i$, and $wv_j$. In the first two cases the
second coordinates $r_i$ and $r_j$ are not adjacent in $W_K$, as
$j\notin\{i-1,i+1\}$ and $W_K$ has no loops. For $u_iu_j$ the first
coordinates are two shadows, or one shadow twice, which are not adjacent
in~$G$; for $wv_j$ the first coordinates $w$ and $v_{\varphi(j)}$ are not
adjacent in~$G$. So $\iota$ is an isomorphism onto the subgraph of
$G\times W_K$ induced by its image.
\end{proof}

\begin{proof}[Proof of Theorem~\ref{thm:main}]
Let $k=2n+1$ and $K=2m+1\ge5$, and let $G$ be $\mu_k$ or $\mu'_k$. If $n\ge m$,
Lemmas~\ref{lem:hom} and~\ref{lem:graphhom} give an induced copy of $G$ in
$G\times W_K$. If $n\le m$, Lemma~\ref{lem:emb} gives an induced copy of
$\mu_K$ in $G\times W_K$. In both cases the copy is not word-representable by
Theorem~\ref{thm:kp}, since $k\ge K\ge5$ in the first case; so $G\times W_K$
is not word-representable by Lemma~\ref{lem:hered}.
\end{proof}

\section{Formal verification}\label{sec:lean}

The results of this paper are formally verified in Lean~4 (version 4.34.1)
with Mathlib~\cite{mathlib} (version v4.34.1, commit \texttt{d13f23b7}).

\paragraph{Statements before proofs.}
After the proof had been found, the agent that found it wrote a Lean file of
definitions and statements, the \emph{statement file}, which contains no
proofs. An independent agent reviewed version~1 of this file against the
sources and tested its definitions with programs of its own; version~2 applies
its corrections and adds the lemmas used in the proof. A second AI system
reviewed version~2, and it was frozen on 27 September 2026 at 22:04 UTC by
recording its SHA-256 hash
\begin{center}\small
\texttt{40c8efd69f8499c9874f07163043c0793c94c8597b7a3902d5aeb0a2ce0c9b58}.
\end{center}
The proofs were written after the freeze, in separate files that import the
statement file, which was not changed.

\paragraph{What is proved.}
Table~\ref{tab:lean} lists the eleven statements of the frozen file; the file
of checks proves each of them. In ASCII notation, the definition of
word-representability and the statement of Theorem~\ref{thm:main} for $\mu$
read as follows (implicit arguments omitted), with \texttt{forall}, \texttt{exists}, \texttt{->},
\verb|<->|, \verb|/\|, \texttt{in}, \texttt{!=} and \texttt{Not} for
$\forall$, $\exists$, $\to$, $\leftrightarrow$, $\wedge$, $\in$, $\ne$ and
$\neg$.
{\footnotesize
\begin{verbatim}
def Alternate (w : List V) (x y : V) : Prop :=
  noRepeat (w.filter (fun z => z == x || z == y)) = true
def Represents (G : SimpleGraph V) (w : List V) : Prop :=
  (forall v : V, v in w) /\
    forall x y : V, x != y -> (G.Adj x y <-> Alternate w x y)
def WordRepresentable (G : SimpleGraph V) : Prop := exists w : List V, Represents G w
def Problem1Mu : Prop := forall n m : Nat, 1 <= n -> 2 <= m ->
  Not (WordRepresentable (tensorProd (mu (2 * n + 1)) (wheel (2 * m + 1))))
\end{verbatim}}
\noindent Here \texttt{noRepeat} says that no two consecutive letters are
equal, and \texttt{mu n} and \texttt{muExt n} stand for
\texttt{mycielskian (cycleGraph n)} and \texttt{extMycielskian (cycleGraph n)}.
The name \texttt{Problem1Mu} comes from the label of
\cite[Problem~5.1]{alsh} in the source of~\cite{alsh}.
Lemmas~\ref{lem:hered}, \ref{lem:graphhom}, \ref{lem:hom}
and~\ref{lem:emb} are proved in Lean as in the text.

\begin{table}[!htb]\centering\footnotesize
\begin{tabular}{@{}>{\raggedright\arraybackslash}p{0.34\textwidth}
 >{\raggedright\arraybackslash}p{0.5\textwidth}
 >{\raggedright\arraybackslash}p{0.1\textwidth}@{}}\toprule
Lean statement & what it says & here\\\midrule
\code{Problem1Mu}, \code{Problem1MuExt} & $\mu_{2n+1}\times W_{2m+1}$, resp.\
 $\mu'_{2n+1}\times W_{2m+1}$, is not word-representable for all $n\ge1$,
 $m\ge2$ & Thm~\ref{thm:main}\\
\code{Problem5_1} & the conjunction of the two & Thm~\ref{thm:main}\\
\code{MycielskiOddCycleNotWR}, \code{ExtMycielskiOddCycleNotWR} & $\mu_{2t+1}$,
 resp.\ $\mu'_{2t+1}$, is not word-representable for all $t\ge1$ &
 Thm~\ref{thm:kp}\\
\code{WordRepresentableOfEmbedding} & a graph that embeds as an induced
 subgraph into a word-representable graph is word-representable &
 Lem.~\ref{lem:hered}\\
\code{EmbeddingOfHom} & a homomorphism $G\to H$ gives an embedding of $G$ into
 $G\times H$ as an induced subgraph & Lem.~\ref{lem:graphhom}\\
\code{Lemma1Mu}, \code{Lemma1MuExt} & a homomorphism $\mu_{2n+1}\to W_{2m+1}$,
 resp.\ $\mu'_{2n+1}\to W_{2m+1}$, for $2\le m\le n$ & Lem.~\ref{lem:hom}\\
\code{Lemma2Mu}, \code{Lemma2MuExt} & an embedding of $\mu_{2m+1}$ as an induced
 subgraph into $\mu_{2n+1}\times W_{2m+1}$, resp.\
 $\mu'_{2n+1}\times W_{2m+1}$, for $1\le n\le m$, $m\ge2$ & Lem.~\ref{lem:emb}\\
\bottomrule
\end{tabular}
\caption{The eleven statements of the frozen statement file (namespace
\texttt{WordRepTensor.\allowbreak Challenge}). Each is proved.}\label{tab:lean}
\end{table}

\paragraph{How Theorem~\ref{thm:kp} is proved in Lean.}
The Lean proof does not follow~\cite{hameed,kp}; it works with words
directly. Let $x\ne y$ be letters of a word, and let $x$ occur first. Then
$x$ and $y$ alternate if and only if every prefix of the word contains at
least as many letters~$x$ as letters~$y$, and at most one more. Now let a word
represent a graph, and orient each edge from the end vertex that occurs
first. Write $\#x$ for the number of letters~$x$ in a prefix. If
$a\to c\to d\to b$ is a directed path and $a$ is adjacent to~$b$, then $a$
occurs before~$b$, so in every prefix each of $\alpha=\#a-\#c$,
$\beta=\#c-\#d$, $\gamma=\#d-\#b$ and $\delta=\#a-\#b$ is $0$ or~$1$.
Hence $\#a-\#d=\alpha+\beta=\delta-\gamma$ and
$\#c-\#b=\beta+\gamma=\delta-\alpha$ are $0$ or~$1$; as $a$ occurs
before~$d$, and $c$ before~$b$, the vertex $a$ is adjacent to~$d$, and $c$
to~$b$. For $\mu_{2t+1}$ and $\mu'_{2t+1}$ with $t\ge2$, consider the seven
vertices
$w,v_i,v_{i+1},v_{i+2},u_i,u_{i+1},u_{i+2}$. A computation checked by the
Lean kernel runs through all orientations of the edges among them ($2^9$ for
$\mu$, $2^{11}$ for $\mu'$) and shows that in every orientation with the
property just proved, the orientation of the five edges between the columns
$i$ and $i+1$ (the edges $v_iv_{i+1}$, $v_iu_{i+1}$, $u_iv_{i+1}$, $wu_i$
and $wu_{i+1}$) and that of the five corresponding edges between the columns
$i+1$ and $i+2$ get different colours, for a fixed colouring of the $2^5$
orientations of five edges with two colours. Going once around the odd
cycle changes the colour an odd number of times, which is impossible. The graphs $\mu_3$
and $\mu'_3$ are equal, and a computation over the $2^{12}$ orientations of
their twelve edges shows that none has the property.

\paragraph{Checks.}
For every statement of Table~\ref{tab:lean}, the axiom audits
(\texttt{\#print axioms}) report only the axioms \texttt{propext},
\texttt{Classical.choice} and \texttt{Quot.sound}. A static scan found no
\texttt{sorry}, \texttt{admit}, \texttt{native\_decide},
\texttt{bv\_decide}, \texttt{implemented\_by}, \texttt{extern},
\texttt{unsafe} or \texttt{opaque}, and no \texttt{axiom} declaration, in
the proof files. The Lean package of this paper, $11$ modules (the
statement file, its tests, eight proof modules and the file of checks), was
built module by module in an empty directory,
reusing only the pinned toolchain and the compiled Mathlib and its
dependencies, and every module was checked with \texttt{leanchecker}, all
with exit code~$0$, twice (on 28 September 2026 between 20:12 and 20:21~UTC).
For the release, the whole Lean project of the
repository was rebuilt from a clean directory and all its modules were
checked again with \texttt{leanchecker} (28 September 2026,
20:48--22:16~UTC). The Lean
files are in version~\repoversion{} of the repository~\cite{repo},
directory \texttt{lean/Research/WordRepTensor/}.

\paragraph{Definitions the reader must trust.}
Apart from Mathlib's graphs, cycle graphs (\texttt{cycleGraph}),
homomorphisms, embeddings and lists, the statements use the definitions of
alternation, representation and word-representability shown above; the
tensor product (\texttt{tensorProd}, built with Mathlib's
\texttt{SimpleGraph.fromRel}); the wheel (\texttt{wheel}, on
\texttt{Sum (Fin n) Unit}); and the Mycielskian and the extended Mycielskian
(\texttt{mycielskian}, \texttt{extMycielskian}, on
\texttt{Sum V (Sum V Unit)},
with top vertices, shadows and root in this order). A file of known-answer
tests checks these definitions on small graphs: adjacencies and edge counts
of cycles, wheels, Mycielskians and tensor products; words that represent
$C_4$, $C_5$, $K_3$ and $K_2\times K_2$; words that do not represent a given
graph, including words that omit a vertex; small instances of the maps of
Lemmas~\ref{lem:graphhom}, \ref{lem:hom} and~\ref{lem:emb}; and a wrong
map that is not an embedding.

\paragraph{Faithfulness and non-claims.}
A word is a list of vertices in which every vertex must occur, as
in~\cite{alsh}; so the empty word represents no graph with a vertex. The Lean
statements concern word-representability itself, not semi-transitive
orientations, so the theorem of~\cite{hkp} is not needed as a hypothesis.
Theorem~\ref{thm:kp} is formalized in its word-representability form; we do
not formalize the characterizations of~\cite{hameed,kp} for arbitrary graphs.
That the Lean definitions agree with those of~\cite{alsh} was checked by the
reviewing agents, not by a machine.

\section*{How this paper was produced}
This work was carried out from 27 to 29 September 2026 (UTC) in a workflow
directed by the author, using AI agents. Agents of Anthropic Claude Code,
with the model Claude Opus~5.5, selected the question, found the proof,
checked it with two separately written programs for small parameters, wrote
the statement file of Section~\ref{sec:lean}, carried out the main
literature search described in the introduction, and drafted this text. An independent
Claude Code agent reviewed version~1 of the statement file against the
sources and checked the definitions with its own programs; its corrections
were applied in version~2. OpenAI Codex (desktop application; models
\texttt{gpt-6-astra} and \texttt{gpt-6-sol}) then checked the
mathematics independently, recomputing the embeddings with its own program,
reviewed version~2 of the statement file and froze it. When Codex's usage
quota ran out, Claude Code wrote the Lean proofs and ran the checks of
Section~\ref{sec:lean}, and an independent Claude Code agent reviewed the
proofs. Another independent Claude Code agent refereed a draft of this
paper. Codex has not reviewed the Lean proofs or this paper; the final
review before release was done by an independent Claude Code agent in its
place.
Throughout, ``we'' refers to this workflow. The author is not a professional
mathematician.
% Author role chosen by the user on 2026-09-28T23:35Z (option A).
The author has read the paper and, using a side-by-side table, compared the
statements of the Lean theorems with the results stated here.
No human mathematician has reviewed this paper. The correctness of the
result rests on the Lean proofs of the statements in Table~\ref{tab:lean}
and on the agreement of those statements with the definitions
of~\cite{alsh}, which was checked by the reviewing agents; the written proofs
are given for the reader. AI systems are not authors of this paper; the
author takes full responsibility for its content.

\medskip
\noindent Thanks are due to Nawaf Alshammari for the question studied here.

{\small
\begin{thebibliography}{99}
\bibitem{alsh}
N.~Alshammari, On the word-representability of tensor product graphs,
arXiv:2609.20881v1, 16 September 2026; we use its numbering.
\bibitem{ckk}
I.~Choi, J.~Kim and M.~Kim, On operations preserving semi-transitive
orientability of graphs, \emph{J.\ Comb.\ Optim.} \textbf{37} (2019),
1351--1366, \url{https://doi.org/10.1007/s10878-018-0358-7};
arXiv:1810.11620.
\bibitem{repo}
H.~Dong, \emph{ai4math-results}, repository,
\url{https://github.com/Horace-Maxwell/ai4math-results}, version~\repoversion{}
(2026); archived at Zenodo under the concept DOI
\url{https://doi.org/10.5281/zenodo.22970254}.
\bibitem{hkp}
M.~M. Halld\'orsson, S.~Kitaev and A.~Pyatkin, Semi-transitive orientations
and word-representable graphs, \emph{Discrete Appl.\ Math.} \textbf{201}
(2016), 164--171, \url{https://doi.org/10.1016/j.dam.2015.07.033}.
\bibitem{hameed}
H.~Hameed, On semi-transitivity of (extended) Mycielski graphs,
\emph{Discrete Appl.\ Math.} \textbf{359} (2024), 83--88,
\url{https://doi.org/10.1016/j.dam.2024.07.028}.
\bibitem{kl}
S.~Kitaev and V.~Lozin, \emph{Words and Graphs}, Monographs in Theoretical
Computer Science, an EATCS Series, Springer, 2015,
\url{https://doi.org/10.1007/978-3-319-25859-1}.
\bibitem{kp}
S.~Kitaev and A.~Pyatkin, A note on Hameed's conjecture on the
semi-transitivity of Mycielski graphs, \emph{Discuss.\ Math.\ Graph Theory}
\textbf{45} (2025), \url{https://doi.org/10.7151/dmgt.2575}; quoted from
arXiv:2408.05066v1 (``A note on semi-transitivity of Mycielski graphs''),
whose numbering we use.
\bibitem{mathlib}
The mathlib Community, The Lean mathematical library, in \emph{Proc.\ 9th ACM
SIGPLAN Int.\ Conf.\ Certified Programs and Proofs} (CPP '20), ACM, 2020,
\url{https://doi.org/10.1145/3372885.3373824}.
\bibitem{scope}
SCOPE-Science, \texttt{SCOPE2026} repository,
\url{https://github.com/SCOPE-Science/SCOPE2026}, accessed 27 September 2026.
\bibitem{tru}
The Omega Institute, \texttt{trureturing} repository,
\url{https://github.com/the-omega-institute/trureturing}, accessed 27
September 2026.
\end{thebibliography}}
\end{document}
